\documentclass[12pt]{article}

\usepackage{amsmath}
\usepackage{german}
\usepackage{graphicx}
\DeclareMathOperator\sgn{sgn}

%\graphicspath{{/home/ritz/Teaching/fcn/fig/}}

%----------------- layout ----------------------------------------

\oddsidemargin=0in
\textwidth=6in
\topmargin=-0.7in
\textheight=9.5in
\baselineskip=16pt
\font\mittel=cmr17
\font\roman=cmr12
\font\klein=cmr9
\font\mgross=cmr10
\newfont{\CAP}{cmcsc10 scaled 900}
\pagestyle{empty}
\parindent=0mm

%----------------- macros ----------------------------------------
\newcommand{\conc}[1]{[{\rm C}]_{\text{#1}}}
\newcommand{\units}[1]{\text{#1}}
\newcommand{\K}{\text{K}}
\newcommand{\Na}{\text{Na}}
\newcommand{\Cl}{\text{Cl}}
\newcommand{\Leak}{\text{leak}}
\newcommand{\water}{\text{H}_2\text{O}}
\newcommand{\tderiv}[1]{\frac{\mathrm{d}#1}{\mathrm{d}t}}
%-----------------------------------------------------------------


\begin{document}

\parbox{2cm}{
\includegraphics[width=1.8cm]{bccnlogo.pdf}
}
\parbox{11cm}{
\begin{center}
\large HUMBOLDT-UNIVERSIT"AT \hskip 0.1 cm ZU \hskip 0.1 cm BERLIN
\vskip 0.1 true cm
\mgross BERNSTEIN CENTER FOR\\ COMPUTATIONAL NEUROSCIENCE
\end{center}
}
\parbox{2cm}
{
\hfill
\includegraphics[width=1.8cm]{hublogo.pdf}
}

\vskip 0.6 true cm

\leftline{\CAP Humboldt-Universit"at zu Berlin 
\hfill Phone: 030/2093-9110}
\leftline{\CAP Philippstr. 13 House 6
              \hfill Fax: 030/2093-6771}
\leftline{\CAP \hfill webpage: http://www.bccn-berlin.de/}


\vskip 0.8 true cm
\centerline{\bf Models of Neural Systems, WS 2010/11}
\centerline{\bf Project 5: Stochastic ion channels}
\centerline{Project presentation and report submission: February, 15th, 2011}

\vskip 0.8 true cm
\noindent {\bf Background}

The Hodgkin-Huxley formalism accurately describes voltage-dependent
conductances of a large number of channels. However, in order to
simulate single channels, we need a model of their different conformational
states and transitions between them. In this project you will
implement a simple model of multi-state ion channels, simulate them
using a Monte-Carlo method, and compare them with conductance equations
proposed by Hodgkin and Huxley.

\vskip 0.8 true cm

{\bf Problems}


\begin{enumerate}
    \item Implement a single delayed-rectifier $\K^+$ channel according to
       the Figure 5.12 in Dayan and Abbott \cite{dayan2001}. Compare
       the obtained channel current to the prediction of Hodgkin-Huxley
       model. Increase the number of channels and see how the
       correspondence between both models changes. Plot the mean
       squared deviation between the current produced by model
       channels and the Hodgkin-Huxley current as a function of N
       over the range 1 $<$ N $<$ 100, matching the amplitude of the
       Hodgkin-Huxley model so that the mean currents are the same.
   \item  Implement the $\Na^+$ channel model from Figure 5.13 of Dayan
       and Abbott \cite{dayan2001}. Compare such a channel model and
       the Hodgkin-Huxley prediction for different numbers of channels.
       Can you see any differences between them? Where do they come
       from?
    \item Simulate the Hodgkin-Huxley neuron model with stochastic
        channels. Plot the distribution of interspike intervals for
        constant current injection. Check how the distribution changes
        with an increasing number of channels. For computational
        efficiency consider implementing improved simulation
        methods (see for example \cite{white2000,chow1996}).
   
   \end{enumerate}
\vspace{1cm}

\begin{thebibliography}{7}
    \bibitem{dayan2001} P Dayan, LF Abbott, ``Theoretical neuroscience'',
        MIT Press 2001
    \bibitem{white2000} JA White, JT Rubinstein, AR Kay, ``Channel
        noise in neurons'', Trends Neurosci. 23: 131-137 (2000)
    \bibitem{chow1996} CC Chow, JA White, ``Spontaneous action
      potentials due to channel fluctuations'', Biophysical Journal
      71: 3013-3021 (1996)
\end{thebibliography}

\vfill
\centerline{\CAP Contact}
\CAP

\begin{tabular}{lll}
\underline{Matthias Guggenmos} & Phone: 450-517131 & Email: matthias.guggenmos(at)bccn-berlin.de \\
Urs Bergmann & Phone: 2093-8924 & Email: urs.bergmann(at)biologie.hu-berlin.de \\
Bartosz Telenczuk & Phone: 2093-8838 & Email: b.telenczuk(at)biologie.hu-berlin.de \\
Richard Kempter & Phone: 2093-8925 & Email: r.kempter(at)biologie.hu-berlin.de \\
\end{tabular}
\end{document}

 
